\documentclass[12pt]{article}

\usepackage[a4paper,margin=1in]{geometry}   % Better page borders
\usepackage{setspace}                       % Line spacing
\usepackage{parskip}                        % Paragraph spacing
\usepackage{abstract}                       % Better abstract formatting

\title{\textbf{The Indian Frontier: Advancing Global Gravitational Wave Astronomy}}
\author{Jitu Moni Kalita}
\date{}

\begin{document}

\maketitle
\doublespacing   % Improves readability

\begin{abstract}
In the early days of astronomy, scientists used only light to study the Universe. With advances in technology and physics, new ways of exploring space have developed. One of the most important discoveries is gravitational waves. This article explains gravitational waves and highlights the contributions of Indian scientists, from early theoretical work to detection efforts and future projects in gravitational-wave astronomy.
\end{abstract}

\section{Introduction}

In 1915, Einstein published his paper on the General Theory of Relativity. In this work, he proposed that gravity is not simply a force of attraction between two massive bodies; rather, it is the bending or curvature of the fabric of space-time. At that time, his idea was considered a radical theory due to the lack of experimental evidence. However, today the Universe provides many confirmations of this curvature of space-time. One such remarkable confirmation is the existence of gravitational waves \cite{overduin}.

The term \textit{gravitational wave} refers to ripples or distortions in the fabric of space-time produced by accelerating massive objects, such as two massive bodies orbiting each other. Nowadays, the detection of gravitational waves has become one of the most important areas in astronomy and astrophysics. It helps scientists study the early Universe and provides valuable insights into phenomena such as dark matter, black holes, and neutron stars. The observation of gravitational waves has opened a new window for scientists to explore and better understand the Universe \cite{jia}. Gravitational waves can be detected using several methods and observatories, such as \textbf{LIGO} (Laser Interferometer Gravitational-Wave Observatory), \textbf{LISA} (Laser Interferometer Space Antenna), and \textbf{PTA} (Pulsar Timing Array).

The contributions of Indian scientists to gravitational-wave observation and understanding have been remarkable. While scientists around the world focused primarily on developing detection methods for gravitational waves, Indian researchers made significant contributions to improving data analysis techniques, particularly by developing methods to enhance the \textit{signal-to-noise ratio} (SNR), which helps extract clearer signals from noisy observational data.

\begin{thebibliography}{9}

\bibitem{overduin}
Overduin, J. (2007). \textit{Gravity as curved spacetime}. In \textit{Exploring Black Holes: Introduction to General Relativity}. Stanford University. Available at: https://einstein.stanford.edu/SPACETIME/spacetime2.html.

\bibitem{jia}
Jia, C. (2023). \textit{Gravitational Waves Through Time: Scientific Significance, Detection Techniques, and Recent Breakthroughs}. Physics and Astronomy Department, University of Texas at Austin.

\end{thebibliography}

\end{document}